\section*{Limitations}
\label{app:limitations}

Four of five generators are open-weight; GPT-5.4 is a non-reproducible Azure snapshot (Appendix~\ref{app:judges}) that also serves as judge and rewrite meta-judge. All primary claims replicate on the open-weight generators independently of it; as a generator, GPT-5.4's self-bias when judged stays within range on other generators except sequence accuracy ($+0.172$ vs.\ $\leq+0.008$), a post-hoc correlational signal we cannot rule out and, being non-reproducible, cannot independently re-verify (Appendix~\ref{app:limitations-extended}).

We use a deterministic programmatic scorer, rather than a human or LLM annotator, as the reference signal, since only it scales to $321{,}648$ evaluations without reintroducing the reliability question under study. A 120-record single-annotator human study (Appendix~\ref{app:prog-judge-validation}) validates it at $92.5$--$98.3\%$ agreement on three of four metrics but only $82.5\%$ on parameter structure, where the scorer penalises schema-valid extra keys that annotators accept ($9$ of $21$ tool-selection disagreements reflect an analogous tool-redundancy gap); this does not drive our findings (ranking shifts $\leq0.09$\,pp under a maximally generous correction), though headline \emph{absolute} numbers still reflect the current scoring rule, so we report both programmatic and human-verdict rankings for best-judge claims. Extended discussion is in Appendix~\ref{app:limitations-extended}.

Records are synthetically generated rather than mined from enterprise traces, since a verified per-record reference -- required by our paired GT/no-GT protocol -- is not obtainable at this scale from proprietary systems; a two-level quality gate and the same human study ($92.7\%$ agreement) validate the pipeline, but domain drift remains an open concern (Appendix~\ref{app:limitations-extended}). Our six main judges are also all general-purpose: adding Prometheus-2~\citep{kim2024prometheus2} as a judge-specialised baseline shows it does not cluster with them ($20$--$30$\,pp lower alignment; Table~\ref{tab:interjudge}) and it is excluded from six-judge statistics (Appendix~\ref{app:limitations-extended}).

The RQ4/RQ6 ablations were each extended to a second (judge, generator) pairing: temperature insensitivity replicates tightly, but the prompt-format advantage only partially replicates and reverses on hard queries, so we treat format as judge/generator-dependent rather than uniformly dominant. Difficulty rewriting is unanimously validated on $58.1\%$ of easy$\to$medium and $93.9\%$ of medium$\to$hard records, so medium should be read as a robustness check rather than a fully calibrated tier. The without-GT ceiling is also partly prompt-dependent: a C2 ablation (Table~\ref{tab:c2-ceiling}) shifts it $\leq{+1.0}$\,pp under an alternative $0.5$-default prompt for the three strongest generators but $+4.1$--$+5.6$\,pp for the two weakest, so task difficulty is the primary ceiling driver only for capable generators (Appendix~\ref{app:limitations-extended}).

Finally, we study judge reliability at evaluation time only and do not test whether these failure modes carry over as a training signal (e.g., a reward model or model-selection gate); we reason through three findings under this framing -- GT-exposure over-anchoring, the without-GT verdict-compression ceiling, and correlated cross-judge failure -- and flag it as concrete future work in Appendix~\ref{app:limitations-extended}.
